% !TeX spellcheck = pt_BR
\documentclass[12pt,portuguese]{article}

%----------------------------------------------------------------------
% PACOTES ESSENCIAIS E TRADUÇÃO
%----------------------------------------------------------------------
\usepackage[brazilian]{babel}     % Suporte à língua portuguesa
\usepackage[utf8]{inputenc}       % Codificação de entrada
\usepackage[T1]{fontenc}          % Codificação da fonte
\usepackage{lmodern}              % Fonte escalável

%----------------------------------------------------------------------
% PACOTES MATEMÁTICOS
%----------------------------------------------------------------------
\usepackage{amsmath, amsthm, amssymb, amstext, mathtools}
\usepackage{mathrsfs}             % Para fontes como \mathscr
\usepackage{dsfont}               % Para o indicador de conjunto \mathds{1}
\usepackage[all]{xy}              % Para diagramas comutativos
\usepackage{imakeidx}             % Criação de Lista de Símbolos / Índice

% Desigualdades sempre com símbolos inclinados
\let\leq\leqslant
\let\le\leqslant
\let\ge\geqslant
\let\geq\geqslant

%----------------------------------------------------------------------
% FORMATAÇÃO, LAYOUT E CORES
%----------------------------------------------------------------------
\usepackage{geometry}
\usepackage{varwidth}
\usepackage{graphicx}
\usepackage[pdftex]{xcolor}
\usepackage{titlesec}
\usepackage{appendix}
\usepackage{enumitem}
\usepackage{float}
\usepackage{booktabs}
\usepackage{array}
\usepackage{colortbl}
\usepackage{fontawesome5}
\usepackage{lucide-icons}
\usepackage{caption}
\usepackage{stackengine}
\captionsetup{width=0.8\textwidth}

% Definição de Cores
\definecolor{teal}{RGB}{15,118,110}
\definecolor{text}{RGB}{28,25,23}
\definecolor{muted}{RGB}{87,83,78}
\definecolor{card}{RGB}{255,255,255}
\definecolor{border}{RGB}{231,229,228}

% Definição de margens e espaçamento
\setlength{\topmargin}{-0.7in}
\setlength{\oddsidemargin}{0in}
\setlength{\textheight}{9.55in}
\setlength{\textwidth}{6.5in}
\setlength{\baselineskip}{12mm}

% Formatação dos títulos de seção
\titleformat{\section}{\Large\bfseries}{\thesection.}{0.5em}{}
\titleformat{\subsection}{\large\bfseries}{\thesubsection.}{0.5em}{}
\titleformat{\subsubsection}{\normalfont\itshape}{\thesubsubsection.}{0.6em}{}

%----------------------------------------------------------------------
% TIKZ E GRÁFICOS
%----------------------------------------------------------------------
\usepackage{tikz}
\usepackage{pgfplots}
\pgfplotsset{compat=1.17}
\usetikzlibrary{positioning,calc,fit,nfold,decorations.pathreplacing,calligraphy,tikzmark,patterns,arrows.meta,bending,hobby}

%----------------------------------------------------------------------
% REFERÊNCIAS, LINKS E METADADOS DO PDF
%----------------------------------------------------------------------
\usepackage[pdftex, colorlinks, linkcolor={red!50!black}, 
citecolor={red}, urlcolor={blue!80!black}]{hyperref}

\hypersetup{
pdftitle={Modos de Convergência de Variáveis Aleatórias},
pdfauthor={L. Cioletti},
pdfsubject={Convergência quase certa, em probabilidade, em média, em média quadrática e em distribuição},
pdfkeywords={convergência quase certa, convergência em probabilidade, convergência em média quadrática, convergência em distribuição, lema de Borel-Cantelli, MAT-UnB}
}

% --- PACOTE E COMANDO PARA A LICENÇA CREATIVE COMMONS ---
\usepackage[type={CC},modifier={by-nc-sa},version={4.0}]{doclicense}
\newcommand\blfootnote[1]{%
\begingroup
\renewcommand\thefootnote{}\footnote{#1}%
\addtocounter{footnote}{-1}%
\endgroup
}

%----------------------------------------------------------------------
% ALIASCNT E AMBIENTES DE TEOREMA (Cleveref)
%----------------------------------------------------------------------
\usepackage{aliascnt}
\theoremstyle{definition}
\newtheorem{theorem}{Teorema}

\newaliascnt{definition}{theorem}
\newtheorem{definition}[definition]{Definição}
\aliascntresetthe{definition}

\newaliascnt{example}{theorem}
\newtheorem{example}[example]{Exemplo}
\aliascntresetthe{example}

\newaliascnt{proposition}{theorem}
\newtheorem{proposition}[proposition]{Proposição}
\aliascntresetthe{proposition}

\newaliascnt{corollary}{theorem}
\newtheorem{corollary}[corollary]{Corolário}
\aliascntresetthe{corollary}

\newaliascnt{lemma}{theorem}
\newtheorem{lemma}[lemma]{Lema}
\aliascntresetthe{lemma}

\newaliascnt{remark}{theorem}
\newtheorem{remark}[remark]{Observação}
\aliascntresetthe{remark}

\newtheoremstyle{plainexercise}
  {3pt}{3pt}{\normalfont}{}{\bfseries}{.}{.5em}
  {\thmnumber{#2}\thmnote{ #3}}
\theoremstyle{plainexercise}
\newtheorem{exercise}{Exercício}

% Versão marcada dos exercícios: a estrela fica pendurada na margem esquerda,
% como nas notas anteriores.
\newtheoremstyle{starredexercise}
  {3pt}{3pt}{\normalfont}{}{\bfseries}{.}{.5em}
  {\llap{\textasteriskcentered\,}\thmnumber{#2}\thmnote{ #3}}
\theoremstyle{starredexercise}
\newtheorem{starredexercise}[exercise]{Exercício}
\theoremstyle{definition}

% Ambientes não numerados usados no texto
\newtheorem*{property}{Propriedade}
\newtheorem*{solution}{Solução}

% Se um enumerate for o primeiro conteúdo de um teorema/exercício (cabeçalho
% ainda pendente, \@inlabel true), o \leavevmode dispara o \everypar do amsthm
% que coloca o cabeçalho na própria linha, antes da lista iniciar.
\makeatletter
\setlist[enumerate]{before={\if@inlabel\leavevmode\fi}}
\makeatother

% Nomes em português para os ambientes
\AtBeginDocument{\renewcommand{\proofname}{\normalfont\textbf{Prova}}}
\renewcommand{\qedsymbol}{$\blacksquare$}

% Cleveref (deve vir depois de hyperref e amsthm)
\usepackage[brazilian,nameinlink]{cleveref}

\crefname{theorem}{teorema}{teoremas}
\Crefname{theorem}{Teorema}{Teoremas}
\crefname{definition}{definição}{definições}
\Crefname{definition}{Definição}{Definições}
\crefname{lemma}{lema}{lemas}
\Crefname{lemma}{Lema}{Lemas}
\crefname{proposition}{proposição}{proposições}
\Crefname{proposition}{Proposição}{Proposições}
\crefname{corollary}{corolário}{corolários}
\Crefname{corollary}{Corolário}{Corolários}
\crefname{example}{exemplo}{exemplos}
\Crefname{example}{Exemplo}{Exemplos}
\crefname{exercise}{exercício}{exercícios}
\Crefname{exercise}{Exercício}{Exercícios}
\crefname{remark}{observação}{observações}
\Crefname{remark}{Observação}{Observações}
\crefname{section}{seção}{seções}
\Crefname{section}{Seção}{Seções}

%----------------------------------------------------------------------
% MACROS ADICIONAIS
%----------------------------------------------------------------------
\let\oldemptyset\emptyset
\let\emptyset\varnothing
\graphicspath{{./}{../arquivos-aux/}}

% Seno na convenção em português
\renewcommand{\sin}{\operatorname{sen}}

%----------------------------------------------------------------------
% DOCUMENTO PRINCIPAL
%----------------------------------------------------------------------
\begin{document}
% --- CABEÇALHO COM O LOGO DA UnB ---
\vspace*{-0.8cm}
\begin{figure}[H]
\includegraphics{logo-mat-unb.pdf}
\end{figure}


% --- TÍTULO, AUTOR E DATA ---
\vspace*{0.3cm}
\begin{center}
{\LARGE\bf Modos de Convergência\\[2pt] de Variáveis Aleatórias}
\vspace{1cm}

{\large \bf L. Cioletti}
\vspace{0.3cm}

{\footnotesize Outubro de 2026}
\end{center}

% --- RESUMO (ABSTRACT) ---
\begin{abstract}
Nestas notas, estudamos as noções de convergência para sequências de variáveis
aleatórias. Começamos recordando os limites superior e inferior de eventos e o
lema de Borel-Cantelli, que são a ferramenta técnica de quase todos os
argumentos. Introduzimos então a convergência quase certa, a convergência em
probabilidade, a convergência em média, a convergência em média quadrática e a
convergência em distribuição, discutindo em cada caso o significado da
definição e o papel dos quantificadores. Mostramos que a convergência em média
quadrática implica as convergências em média, em probabilidade e em
distribuição, que a convergência em probabilidade implica a convergência em
distribuição e que a convergência quase certa implica a convergência em
probabilidade. Para cada implicação, exibimos um contraexemplo que mostra por
que a recíproca falha. Estudamos também o critério de Borel-Cantelli para a
convergência quase certa, com uma discussão detalhada da definição por
$\varepsilon$ e $n_0$, e o teorema da convergência dominada, que mostra como a
convergência quase certa de uma sequência dominada implica todas as demais. O
texto inclui exemplos e uma lista de exercícios ao final. A exposição segue de
perto as Seções 1 e 2 do Capítulo 1 e a Seção 25 do Capítulo 5 de
\cite{MR2893652}, as Seções 8.1, 8.2 e 8.5 do Capítulo 8 de
\cite{GrimmettWelsh2014} e a Seção 2.3 de \cite{MR3930614}.
\end{abstract}

% --- LICENÇA NO RODAPÉ ---
\blfootnote{%
\doclicenseImage[imagewidth=2.6em]\ 
\begingroup
\hypersetup{urlcolor=darkgray}
\href{\doclicenseURL}{%
\doclicenseName - Attribution-NonCommercial-ShareAlike 4.0 International%
}%
\endgroup
}

%----------------------------------------------------------------------
% CONTEÚDO
%----------------------------------------------------------------------
\section{Introdução}\label{sec:introducao}

Nas notas anteriores, estudamos uma variável aleatória por vez. Fixado o
espaço de probabilidade $(\Omega,\mathcal{F},\mathbb{P})$, descrevemos a
distribuição de $X$, calculamos sua esperança, sua variância e seus momentos e
usamos essas grandezas para estimar a probabilidade de eventos associados a
$X$. A desigualdade de Chebyshev, por exemplo, transforma informação sobre a
variância em uma cota para $\mathbb{P}(|X-\mathbb{E}[X]|\geqslant t)$. Em todas
essas situações, o objeto de interesse era um único experimento aleatório.

Nesta nota, o objeto de interesse passa a ser uma sequência
$X_1,X_2,X_3,\ldots$ de variáveis aleatórias definidas no mesmo espaço de
probabilidade. Sequências aparecem naturalmente em vários problemas. Ao lançar
uma moeda repetidamente, podemos acompanhar a frequência relativa de caras nos
$n$ primeiros lançamentos e perguntar se ela se estabiliza. Ao repetir um
experimento para estimar uma média desconhecida, como na integração de Monte
Carlo, calculamos médias de um número crescente de observações e queremos saber
se elas se aproximam do valor procurado. Ao considerar o máximo
$M_n=\max\{U_1,\ldots,U_n\}$ de $n$ variáveis uniformes em $(0,1)$, vemos que
$M_n$ aumenta com $n$ e parece se aproximar de $1$. Em cada um desses exemplos,
temos uma sequência de números aleatórios e uma pergunta do tipo: ela se
aproxima de algum limite? Em que sentido? Com que garantia?

A expressão ``em que sentido'' é o ponto central da nota, porque há várias
respostas razoáveis e elas não são equivalentes. A resposta mais direta é
exigir que, para cada ponto $\omega$ do espaço amostral, a sequência de números
reais $X_n(\omega)$ convirja para $X(\omega)$ no sentido usual da análise. Essa
é a convergência quase certa. Ela é mais forte que a convergência em
probabilidade e que a convergência em distribuição, mas não é comparável com as
convergências em média e em média quadrática, como veremos na
\Cref{sec:medias}.
Há, porém, duas dificuldades. A primeira é que a convergência pode falhar em
um conjunto de pontos, e esse conjunto precisa ser pequeno do ponto de vista da
medida de probabilidade, não do ponto de vista da cardinalidade. A segunda é
que, em muitos problemas concretos, não há como controlar o comportamento
ponto a ponto de toda a sequência: o que sabemos estimar é a probabilidade de
que $X_n$ se afaste de $X$ em uma etapa fixada $n$. Isso sugere uma noção mais
fraca, na qual se pede que a probabilidade de um desvio de tamanho
$\varepsilon$ tenda a zero quando $n$ cresce. Essa é a convergência em
probabilidade.

Entre essas duas noções há um espaço intermediário, ocupado pelas
convergências definidas por médias dos erros. Se o erro $X_n-X$ tem média
pequena, então, por Markov, a probabilidade de que ele seja grande também é
pequena. Assim, exigir que $\mathbb{E}|X_n-X|$ tenda a zero, a convergência em
média, ou que $\mathbb{E}[(X_n-X)^2]$ tenda a zero, a convergência em média
quadrática, produz noções mais fortes que a convergência em probabilidade. A
vantagem dessas noções é que elas se reduzem a cálculos de esperança, com os
quais já estamos habituados. A desvantagem é que a média pode ser estragada por
poucos pontos onde o erro é enorme, mesmo que a sequência convirja em quase
todos os pontos.

Há ainda uma quarta possibilidade, que abandona a comparação entre os valores
de $X_n$ e os de $X$ e compara apenas as distribuições. Dizemos que $X_n$
converge em distribuição para $X$ quando as funções de distribuição
$F_{X_n}(x)=\mathbb{P}(X_n\leqslant x)$ convergem para $F_X(x)$ em todo ponto
$x$ em que $F_X$ é contínua. Essa noção é a mais fraca de todas e é a que
aparece no teorema central do limite, quando afirmamos que as somas
normalizadas têm distribuição aproximadamente normal, sem que isso signifique
que elas se aproximem de uma variável normal específica. A convergência em
distribuição também é a única das noções que faz sentido quando as variáveis
$X_n$ vivem em espaços de probabilidade diferentes, ou quando a variável
limite não está definida no mesmo espaço que as demais.

O objetivo da nota é apresentar as cinco noções, compará-las e fornecer
ferramentas para verificar convergência em problemas concretos. A
\Cref{sec:preliminares} recorda os resultados sobre limites de eventos e o lema
de Borel-Cantelli que serão usados repetidamente; essas ferramentas já foram
demonstradas na nota \emph{Desigualdades de Markov e Chebyshev e o lema de
Borel-Cantelli}, e aqui apenas as enunciamos e comentamos. Na
\Cref{sec:quase-certa}, tratamos da convergência quase certa e discutimos com
cuidado a definição por $\varepsilon$ e $n_0$, que costuma ser a primeira
dificuldade real do assunto. A \Cref{sec:probabilidade} introduz a convergência
em probabilidade e mostra que ela é consequência da convergência quase certa,
mas não é equivalente a ela. A \Cref{sec:medias} discute as convergências em
média e em média quadrática e as comparações entre elas. A
\Cref{sec:distribuicao} introduz a convergência em distribuição e prova que a
convergência em probabilidade a implica. A \Cref{sec:quadro} reúne todas as
implicações em um diagrama, resume as definições em uma tabela e apresenta o
teorema da convergência dominada, que mostra como a convergência quase certa,
acompanhada de uma cota integrável, recupera todas as outras noções.

Um fio condutor percorre o texto. Quase toda prova de convergência quase certa
passa por uma estimativa da forma
\[
\sum_{n=1}^{\infty}\mathbb{P}\bigl(|X_n-X|\geqslant\varepsilon\bigr)<\infty,
\]
seguida do primeiro lema de Borel-Cantelli. O motivo é que essa série mede o
custo total, ao longo de toda a sequência, dos desvios de tamanho
$\varepsilon$. Se o custo total é finito, então, para quase todo $\omega$,
apenas finitos desvios ocorrem, e a sequência converge. Esse é o conteúdo da
\Cref{prop:criterio-bc}, que será usado na nota sobre a lei forte dos grandes
números. Vale registrar desde já que convergência em probabilidade fornece
apenas cotas que tendem a zero, sem que a soma convirja, e é exatamente essa
diferença que separa as duas noções.

A exposição segue de perto as Seções 1 e 2 do Capítulo 1 e a Seção 25 do
Capítulo 5 de \cite{MR2893652}, as Seções 8.1, 8.2 e 8.5 do Capítulo 8 de
\cite{GrimmettWelsh2014} e a Seção 2.3 de \cite{MR3930614}. Os exemplos foram
escolhidos para que cada implicação e cada recíproca falsa tenham pelo menos um
exemplo concreto associado.

%----------------------------------------------------------------------
\section{Preliminares: limites de eventos e Borel-Cantelli}\label{sec:preliminares}

Nesta seção, recordamos os resultados sobre sequências de eventos que serão
usados em quase todas as demonstrações. As provas estão na nota
\emph{Desigualdades de Markov e Chebyshev e o lema de Borel-Cantelli}, onde
também se encontra a discussão completa dos limites superior e inferior. Aqui
repetimos os enunciados para que o texto possa ser lido de forma contínua e
acrescentamos apenas os comentários necessários ao uso que faremos deles.

\begin{definition}[Limite superior de uma sequência de eventos]\label{def:limsup}
Seja $(A_n)_{n\geqslant1}$ uma sequência de eventos. O \emph{limite superior}
de $(A_n)$ é o evento
\[
\limsup_{n\to\infty}A_n
=
\bigcap_{n=1}^{\infty}\bigcup_{k=n}^{\infty}A_k.
\]
Escrevemos também $\{A_n\text{ infinitas vezes}\}$ ou $\{A_n\ \text{i.v.}\}$
para esse evento.
\end{definition}

Um ponto $\omega$ pertence a $\bigcup_{k\geqslant n}A_k$ quando pertence a pelo
menos um $A_k$ com $k\geqslant n$. Dizer que $\omega$ está na interseção sobre
todos os $n$ significa, portanto, que para cada $n$ existe um índice
$k\geqslant n$ com $\omega\in A_k$, isto é, que os índices $k$ para os quais
$\omega\in A_k$ não são limitados. Em palavras,
\[
\limsup_{n\to\infty}A_n
=
\{\omega\in\Omega:\omega\in A_n\text{ para infinitos índices }n\}.
\]

O limite inferior tem definição dual e corresponde a ``a partir de algum
momento'':
\[
\liminf_{n\to\infty}A_n
=
\bigcup_{n=1}^{\infty}\bigcap_{k=n}^{\infty}A_k
=
\{\omega\in\Omega:\omega\in A_n\text{ para todo }n\text{ suficientemente grande}\}.
\]
Temos $\liminf_n A_n\subseteq\limsup_n A_n$, e o complementar troca um limite
pelo outro:
\[
\Bigl(\limsup_{n\to\infty}A_n\Bigr)^c
=
\liminf_{n\to\infty}A_n^c.
\]

O instrumento que liga esses eventos às suas probabilidades é a continuidade da
medida. Se $(B_n)$ é uma sequência decrescente de eventos e
$B=\bigcap_n B_n$, então
\begin{equation}\label{eq:continuidade}
\mathbb{P}(B)=\lim_{n\to\infty}\mathbb{P}(B_n).
\end{equation}
Aplicando \eqref{eq:continuidade} à sequência decrescente
$C_n=\bigcup_{k\geqslant n}A_k$, cuja interseção é $\limsup_n A_n$, obtemos
\begin{equation}\label{eq:limsup-continuidade}
\mathbb{P}\Bigl(\limsup_{n\to\infty}A_n\Bigr)
=
\lim_{n\to\infty}\mathbb{P}\Bigl(\bigcup_{k=n}^{\infty}A_k\Bigr).
\end{equation}
A desigualdade de Boole, por sua vez, permite estimar a probabilidade de uma
união enumerável:
\begin{equation}\label{eq:boole}
\mathbb{P}\Bigl(\bigcup_{n=1}^{\infty}A_n\Bigr)
\leqslant
\sum_{n=1}^{\infty}\mathbb{P}(A_n).
\end{equation}
Com \eqref{eq:limsup-continuidade} e \eqref{eq:boole}, obtemos a primeira parte
do lema de Borel-Cantelli.

\begin{theorem}[Primeiro lema de Borel-Cantelli]\label{thm:bc1}
Seja $(A_n)_{n\geqslant1}$ uma sequência de eventos tal que
$\sum_{n\geqslant1}\mathbb{P}(A_n)<\infty$. Então
\[
\mathbb{P}\Bigl(\limsup_{n\to\infty}A_n\Bigr)=0.
\]
Em outras palavras, com probabilidade $1$ ocorrem apenas finitos dos eventos
$A_n$.
\end{theorem}

\begin{proof}
Como $C_n=\bigcup_{k\geqslant n}A_k$ é uma sequência decrescente cuja
interseção é $\limsup_n A_n$, a continuidade da medida
\eqref{eq:continuidade} fornece
\[
\mathbb{P}\Bigl(\limsup_{n\to\infty}A_n\Bigr)
=
\lim_{n\to\infty}\mathbb{P}(C_n).
\]
Por \eqref{eq:boole},
$\mathbb{P}(C_n)\leqslant\sum_{k\geqslant n}\mathbb{P}(A_k)$. A cauda de uma
série convergente tende a zero, logo $\mathbb{P}(C_n)\to0$ e o resultado segue.
\end{proof}

A prova não usa independência: o primeiro lema vale para qualquer sequência de
eventos. O segundo lema trata da recíproca, que exige independência.

\begin{theorem}[Segundo lema de Borel-Cantelli]\label{thm:bc2}
Seja $(A_n)_{n\geqslant1}$ uma sequência de eventos independentes com
$\sum_{n\geqslant1}\mathbb{P}(A_n)=\infty$. Então
\[
\mathbb{P}\Bigl(\limsup_{n\to\infty}A_n\Bigr)=1.
\]
\end{theorem}

A prova, que está na nota citada, usa a desigualdade $1-x\leqslant e^{-x}$ para
estimar $\mathbb{P}\bigl(\bigcap_{n=M}^{N}A_n^c\bigr)$ e depois a continuidade
da medida. Não repetiremos o argumento porque ele não será usado diretamente
aqui; o segundo lema aparecerá apenas em exemplos e na discussão da recíproca
do critério de Borel-Cantelli.

Encerramos a seção com uma observação sobre o uso desses resultados. O primeiro
lema transforma uma informação sobre a \emph{soma} das probabilidades em uma
informação sobre o \emph{conjunto} dos pontos onde os eventos ocorrem
infinitas vezes. Essa passagem é a única maneira elementar que temos de obter
convergência quase certa a partir de estimativas de probabilidade, e é por isso
que a série $\sum_n\mathbb{P}(A_n)$ aparecerá tantas vezes no que segue.

%----------------------------------------------------------------------
\section{Convergência quase certa}\label{sec:quase-certa}

\subsection{A definição}\label{subsec:qc-definicao}

A primeira noção de convergência que estudaremos é a mais próxima da
convergência de sequências de números reais. Fixado $\omega$, a sequência
$X_1(\omega),X_2(\omega),\ldots$ é uma sequência de números reais, e podemos
perguntar se ela converge no sentido usual. A convergência quase certa pede que
isso ocorra para quase todo $\omega$.

\begin{definition}[Convergência quase certa]\label{def:qc}
Sejam $X,X_1,X_2,\ldots$ variáveis aleatórias definidas em um mesmo espaço de
probabilidade $(\Omega,\mathcal{F},\mathbb{P})$. Dizemos que $(X_n)$
\emph{converge quase certamente} para $X$ se
\[
\mathbb{P}\Bigl(\bigl\{\omega\in\Omega:
\lim_{n\to\infty}X_n(\omega)=X(\omega)\bigr\}\Bigr)=1.
\]
Nesse caso, escrevemos $X_n\xrightarrow{\text{q.c.}}X$ ou, simplesmente,
$X_n\to X$ quase certamente.
\end{definition}

O conjunto que aparece na definição é o conjunto dos pontos onde a sequência
numérica converge. Denotaremos esse conjunto por
\[
C\equiv\Bigl\{\omega\in\Omega:\lim_{n\to\infty}X_n(\omega)=X(\omega)\Bigr\}.
\]
A definição pede que $C$ tenha probabilidade $1$, ou seja, que seu complementar
tenha probabilidade zero. Os pontos de $C^c$ são os pontos onde a sequência não
converge, e a exigência é que eles formem um conjunto pequeno do ponto de vista
da medida $\mathbb{P}$. Note que $C^c$ pode ser não vazio e pode até ser não
enumerável; o que a definição controla é a sua probabilidade.

Antes de usar a definição, precisamos saber que $C$ é um evento, isto é, que
$C\in\mathcal{F}$. A próxima proposição resolve essa questão e, ao mesmo tempo,
traduz a definição em uma forma que será usada em todas as demonstrações.

\subsection{O conjunto de convergência em termos de eventos}\label{subsec:qc-conjunto}

\begin{proposition}[Caracterização da convergência quase certa]\label{prop:qc-caracterizacao}
Sejam $X,X_1,X_2,\ldots$ variáveis aleatórias. Então
\begin{equation}\label{eq:qc-conjunto}
C
=
\bigcap_{m=1}^{\infty}\bigcup_{N=1}^{\infty}\bigcap_{n=N}^{\infty}
\Bigl\{|X_n-X|<\frac{1}{m}\Bigr\},
\end{equation}
e, equivalentemente,
\begin{equation}\label{eq:qc-limsup}
C^c
=
\bigcup_{m=1}^{\infty}\limsup_{n\to\infty}
\Bigl\{|X_n-X|\geqslant\frac{1}{m}\Bigr\}.
\end{equation}
Em particular, $C\in\mathcal{F}$ e
\[
X_n\xrightarrow{\text{q.c.}}X
\qquad\Longleftrightarrow\qquad
\mathbb{P}\Bigl(\limsup_{n\to\infty}\bigl\{|X_n-X|\geqslant\tfrac{1}{m}\bigr\}\Bigr)=0
\quad\text{para todo }m\geqslant1.
\]
\end{proposition}

\begin{proof}
Provamos primeiro \eqref{eq:qc-conjunto}. Seja $\omega\in C$ e fixe $m\geqslant1$.
Tomando $\varepsilon=1/m$ na definição de convergência, existe $N$ tal que
$|X_n(\omega)-X(\omega)|<1/m$ para todo $n\geqslant N$. Logo $\omega$ pertence a
$\bigcap_{n\geqslant N}\{|X_n-X|<1/m\}$ e, portanto, à união sobre $N$ e à
interseção sobre $m$.

Reciprocamente, suponha que $\omega$ pertence ao conjunto do lado direito de
\eqref{eq:qc-conjunto} e seja $\varepsilon>0$. Escolha $m$ com $1/m<\varepsilon$.
Por hipótese, existe $N$ tal que
$|X_n(\omega)-X(\omega)|<1/m<\varepsilon$ para todo $n\geqslant N$. Como
$\varepsilon>0$ é arbitrário, $\omega\in C$. Isso prova \eqref{eq:qc-conjunto}.

A identidade \eqref{eq:qc-limsup} é a versão complementar de
\eqref{eq:qc-conjunto}. De fato, com $A_{n,m}=\{|X_n-X|<1/m\}$, o complementar
de $\bigcup_N\bigcap_{n\geqslant N}A_{n,m}$ é
\[
\bigcap_N\bigcup_{n\geqslant N}A_{n,m}^c
=
\limsup_{n\to\infty}A_{n,m}^c.
\]
Como $A_{n,m}^c=\{|X_n-X|\geqslant1/m\}$ e o
complementar de uma interseção enumerável é a união dos complementares,
obtemos \eqref{eq:qc-limsup}.

A última afirmação segue de \eqref{eq:qc-limsup} e da subaditividade enumerável
da probabilidade: uma união enumerável de eventos tem probabilidade zero se, e
somente se, cada um deles tem probabilidade zero. Finalmente, como cada
$\{|X_n-X|\geqslant1/m\}$ é um evento, o mesmo vale para o limite superior, para
a união em $m$ e para $C$. Logo $C\in\mathcal{F}$.
\end{proof}

A proposição anterior é a chave de tudo o que segue, e vale a pena reler o seu
enunciado em palavras. Ela diz que a convergência quase certa pode ser testada
tolerância por tolerância: para cada $m$, consideramos o evento em que o erro
$|X_n-X|$ é maior ou igual a $1/m$ em infinitos índices; se todos esses eventos
têm probabilidade zero, então a convergência vale quase certamente. O ponto
delicado é que a interseção em $m$ é enumerável, e é isso que permite usar a
subaditividade da probabilidade. Se tivéssemos que intersectar sobre todos os
$\varepsilon>0$ reais, o argumento não funcionaria, porque uma interseção não
enumerável de eventos de probabilidade $1$ pode ter probabilidade menor que
$1$. A escolha da sequência $1/m$ resolve o problema porque ela é decrescente e
tende a zero.

\subsection{A definição por $\varepsilon$ e $n_0$}\label{subsec:qc-epsilon}

Nesta subseção, discutimos a definição com o cuidado que ela merece, porque é
aqui que aparece a primeira dificuldade real do assunto. A dificuldade não está
na matemática, mas na ordem dos quantificadores.

Em análise, dizemos que $x_n\to x$ quando
\begin{equation}\label{eq:analise}
\text{para todo }\varepsilon>0\text{ existe }N\in\mathbb{N}
\text{ tal que }|x_n-x|<\varepsilon\text{ para todo }n\geqslant N.
\end{equation}
A negação dessa afirmação é
\begin{equation}\label{eq:analise-neg}
\text{existe }\varepsilon>0\text{ tal que, para todo }N\in\mathbb{N},
\text{ existe }n\geqslant N\text{ com }|x_n-x|\geqslant\varepsilon.
\end{equation}
O ponto importante é que, em \eqref{eq:analise}, o número $N$ pode depender de
$\varepsilon$, e em \eqref{eq:analise-neg} o índice $n$ pode depender de $N$.
Ao passar para variáveis aleatórias, acrescentamos uma dependência adicional:
$N$ pode depender de $\omega$. A definição de convergência quase certa não pede
que exista um único $N$ que sirva para todos os pontos; ela pede que, para cada
ponto de um conjunto de probabilidade $1$, exista um $N$ próprio daquele ponto.

Para ver como essa dependência se traduz em eventos, fixe $\varepsilon>0$ e
considere o conjunto dos pontos para os quais a desigualdade
$|X_n-X|<\varepsilon$ vale a partir de algum índice:
\[
B_\varepsilon
\equiv
\bigcup_{N=1}^{\infty}\bigcap_{n=N}^{\infty}
\bigl\{|X_n-X|<\varepsilon\bigr\}.
\]
Um ponto $\omega$ pertence a $B_\varepsilon$ se, e somente se, existe $N$ tal
que $|X_n(\omega)-X(\omega)|<\varepsilon$ para todo $n\geqslant N$. Isso é
exatamente a afirmação ``a tolerância $\varepsilon$ é respeitada a partir de um
certo momento'', com o momento dependendo de $\omega$. O evento $B_\varepsilon$
é, portanto, a versão probabilística da condição que aparece em
\eqref{eq:analise} para uma tolerância fixada.

A convergência em $\omega$ exige que essa condição valha para todas as
tolerâncias ao mesmo tempo. Assim,
\[
C=\bigcap_{\varepsilon>0}B_\varepsilon.
\]
Essa interseção não é enumerável, mas a monotonicidade de $B_\varepsilon$ em
$\varepsilon$ permite substituí-la por uma interseção enumerável: se
$B_{1/m}$ ocorre para todo $m$, então $B_\varepsilon$ ocorre para todo
$\varepsilon>0$, porque basta escolher $m$ com $1/m<\varepsilon$ e usar
$B_{1/m}\subseteq B_\varepsilon$. Logo
\[
C=\bigcap_{m=1}^{\infty}B_{1/m},
\]
que é a identidade \eqref{eq:qc-conjunto}.

Resta entender por que a interseção enumerável preserva a probabilidade $1$.
Suponha que $\mathbb{P}(B_{1/m})=1$ para todo $m$. Então
$\mathbb{P}(B_{1/m}^c)=0$ para todo $m$ e, pela subaditividade enumerável,
\[
\mathbb{P}\Bigl(\bigcup_{m=1}^{\infty}B_{1/m}^c\Bigr)
\leqslant
\sum_{m=1}^{\infty}\mathbb{P}(B_{1/m}^c)=0.
\]
Tomando complementares, $\mathbb{P}(C)=1$. Esse é o único lugar em que a
enumerabilidade é usada, e é também o lugar em que muitos argumentos informais
erram. Não é verdade que uma interseção arbitrária de eventos de probabilidade
$1$ tenha probabilidade $1$; é verdade para interseções enumeráveis.

A negação da convergência quase certa, que será usada para construir
contraexemplos, tem a seguinte forma. O ponto $\omega$ não pertence a $C$ se, e
somente se, existe $m\geqslant1$ tal que, para todo $N$, existe $n\geqslant N$
com $|X_n(\omega)-X(\omega)|\geqslant1/m$. Em símbolos,
\[
\omega\in C^c
\qquad\Longleftrightarrow\qquad
\exists\,m\geqslant1\ \ \forall N\geqslant1\ \ \exists\,n\geqslant N:
\quad |X_n(\omega)-X(\omega)|\geqslant\frac{1}{m}.
\]
Essa é a versão de \eqref{eq:analise-neg} com a dependência em $\omega$
explicitada. A diferença entre
\[
\forall m\ \exists N\ \forall n\geqslant N
\qquad\text{e}\qquad
\exists m\ \forall N\ \exists n\geqslant N
\]
é toda a diferença entre convergência e não convergência, e a segunda forma é
a que devemos verificar quando queremos mostrar que uma sequência \emph{não}
converge em um ponto. Na segunda forma, a tolerância que testemunha a falha é
uma só, escolhida uma vez; ela não precisa funcionar para todos os $m$.

Vale registrar uma consequência prática da discussão. Para provar que
$X_n\to X$ quase certamente, é suficiente exibir, para cada $m$, um evento
$B_m$ de probabilidade $1$ tal que, em $B_m$, a desigualdade
$|X_n-X|<1/m$ valha para todo $n$ suficientemente grande. Em outras palavras,
podemos trabalhar com uma tolerância de cada vez, desde que o façamos para
todas as tolerâncias da sequência $1/m$ e que os eventos excepcionais sejam
enumeráveis.

Essa leitura também mostra o que a convergência em probabilidade não fornece.
Ela dá, para cada tolerância fixada, a informação de que a probabilidade do
evento ruim é pequena em cada etapa, mas não diz nada sobre a organização dos
índices ruins ao longo do tempo. O exemplo da sequência que varre o intervalo,
na \Cref{subsec:typewriter}, mostrará que os conjuntos ruins podem se deslocar
indefinidamente, de modo que cada ponto seja atingido infinitas vezes ainda que
a probabilidade de cada conjunto tenda a zero.

\subsection{Primeiros exemplos}\label{subsec:qc-exemplos}

\begin{example}[Convergência quase certa sem ferramentas]\label{ex:qc-translacao}
Seja $X$ uma variável aleatória qualquer e defina $X_n=X+1/n$. Então
$X_n\to X$ quase certamente. De fato, para todo $\omega$,
\[
|X_n(\omega)-X(\omega)|=\frac{1}{n}\longrightarrow0,
\]
de modo que $C=\Omega$ e $\mathbb{P}(C)=1$. O exemplo é trivial, mas serve para
fixar a notação: a convergência quase certa é uma afirmação sobre a sequência
numérica obtida em cada ponto, e não sobre a distribuição das variáveis.
\end{example}

\begin{example}[O máximo de variáveis uniformes]\label{ex:qc-maximo}
Sejam $U_1,U_2,\ldots$ variáveis aleatórias independentes com distribuição
uniforme em $(0,1)$ e seja $M_n=\max\{U_1,\ldots,U_n\}$. Então
$M_n\to1$ quase certamente. Para $\varepsilon\in(0,1)$, como os
eventos $\{U_k\leqslant1-\varepsilon\}$ são independentes e cada um tem
probabilidade $1-\varepsilon$,
\[
\mathbb{P}\bigl(|M_n-1|\geqslant\varepsilon\bigr)
=
\mathbb{P}\bigl(M_n\leqslant1-\varepsilon\bigr)
=
(1-\varepsilon)^n.
\]
Para cada $m\geqslant2$, tome $\varepsilon=1/m$. A série geométrica
$\sum_n(1-1/m)^n$ converge, logo o primeiro lema de Borel-Cantelli garante que
\[
\mathbb{P}\bigl(|M_n-1|\geqslant\tfrac{1}{m}\ \text{infinitas vezes}\bigr)=0.
\]
Como isso vale para cada $m$ e a união sobre $m$ é enumerável, o conjunto em que
algum desses eventos ocorre tem probabilidade zero. Pela
\Cref{prop:qc-caracterizacao}, $M_n\to1$ quase certamente.
\end{example}

O exemplo anterior antecipa o critério da próxima subseção: a convergência
quase certa foi obtida a partir da somabilidade das probabilidades dos desvios.
O cálculo de $\mathbb{P}(M_n\leqslant1-\varepsilon)$ é elementar porque as
variáveis são independentes, e a somabilidade decorre de a probabilidade decair
geometricamente. Nos problemas de médias, a probabilidade de um desvio decai
apenas como $1/n$, e a série diverge; será preciso um argumento mais fino, que
é o conteúdo da nota sobre a lei forte.

\subsection{O critério de Borel-Cantelli}\label{subsec:qc-criterio}

O resultado a seguir é a ferramenta padrão para provar convergência quase certa.
Ele transforma uma estimativa sobre cada $n$ em uma conclusão sobre o
comportamento de toda a sequência.

\begin{proposition}[Critério de Borel-Cantelli]\label{prop:criterio-bc}
Sejam $X,X_1,X_2,\ldots$ variáveis aleatórias. Suponha que, para todo
$\varepsilon>0$,
\begin{equation}\label{eq:criterio-bc}
\sum_{n=1}^{\infty}\mathbb{P}\bigl(|X_n-X|\geqslant\varepsilon\bigr)<\infty.
\end{equation}
Então $X_n\to X$ quase certamente.
\end{proposition}

\begin{proof}
Fixe $m\geqslant1$ e defina $A_{n,m}=\{|X_n-X|\geqslant1/m\}$. A hipótese,
aplicada a $\varepsilon=1/m$, diz que $\sum_n\mathbb{P}(A_{n,m})<\infty$. Pelo
\Cref{thm:bc1},
\[
\mathbb{P}\Bigl(\limsup_{n\to\infty}A_{n,m}\Bigr)=0.
\]
Como isso vale para cada $m$, a união enumerável
$\bigcup_m\limsup_n A_{n,m}$ também tem probabilidade zero. Pela
\Cref{prop:qc-caracterizacao}, esse conjunto é exatamente $C^c$. Logo
$\mathbb{P}(C)=1$, isto é, $X_n\to X$ quase certamente.
\end{proof}

A leitura da prova esclarece o papel de cada hipótese. O primeiro lema de
Borel-Cantelli é aplicado a uma tolerância fixa; a passagem para todas as
tolerâncias é feita pela união enumerável em $m$; e a identificação do conjunto
excepcional com $C^c$ é feita pela \Cref{prop:qc-caracterizacao}. Nenhuma
hipótese de independência é usada, o que é coerente com o fato de o primeiro
lema não exigir independência.

Vale comparar \eqref{eq:criterio-bc} com a definição de convergência em
probabilidade, que estudaremos na próxima seção. Aquela definição pede apenas
que cada termo $\mathbb{P}(|X_n-X|\geqslant\varepsilon)$ tenda a zero. A
hipótese \eqref{eq:criterio-bc} pede que a soma desses termos seja finita, o
que é estritamente mais forte. A diferença é a mesma que existe entre dizer que
uma sequência de números não negativos tende a zero e dizer que ela é somável:
a segunda afirmação controla o conjunto dos índices, e não apenas o termo
genérico.

Na prática, raramente verificamos \eqref{eq:criterio-bc} para todos os
$\varepsilon$. O corolário seguinte mostra que basta uma sequência de
tolerâncias que tende a zero.

\begin{corollary}[Critério com tolerâncias variáveis]\label{cor:criterio-bc-epsn}
Sejam $X,X_1,X_2,\ldots$ variáveis aleatórias e seja $(\varepsilon_n)$ uma
sequência de números positivos com $\varepsilon_n\to0$. Se
\[
\sum_{n=1}^{\infty}\mathbb{P}\bigl(|X_n-X|\geqslant\varepsilon_n\bigr)<\infty,
\]
então $X_n\to X$ quase certamente.
\end{corollary}

\begin{proof}
Fixe $\varepsilon>0$. Como $\varepsilon_n\to0$, existe $n_0$ tal que
$\varepsilon_n<\varepsilon$ para todo $n\geqslant n_0$. Para esses índices,
\[
\bigl\{|X_n-X|\geqslant\varepsilon\bigr\}
\subseteq
\bigl\{|X_n-X|\geqslant\varepsilon_n\bigr\},
\]
logo
\[
\sum_{n=n_0}^{\infty}\mathbb{P}\bigl(|X_n-X|\geqslant\varepsilon\bigr)
\leqslant
\sum_{n=n_0}^{\infty}\mathbb{P}\bigl(|X_n-X|\geqslant\varepsilon_n\bigr)
<\infty.
\]
Os primeiros $n_0-1$ termos formam uma soma finita e cada um deles é no máximo
$1$, logo a série completa $\sum_{n=1}^{\infty}$ também converge. A
\Cref{prop:criterio-bc} conclui a prova.
\end{proof}

O corolário é a versão que aparece nas aplicações. Na prova da lei forte por
truncamento, por exemplo, estima-se a probabilidade de um desvio de tamanho
$\varepsilon n$, e não de tamanho $\varepsilon$ fixo; a tolerância cresce com
$n$, e o corolário permite concluir a convergência quase certa.

\subsection{O critério não é necessário}\label{subsec:qc-nao-necessario}

A hipótese \eqref{eq:criterio-bc} é suficiente, mas não é necessária. O exemplo
a seguir mostra uma sequência que converge quase certamente e para a qual a
soma das probabilidades dos desvios diverge.

\begin{example}[Convergência quase certa com série divergente]\label{ex:qc-serie-divergente}
Seja $\Omega=(0,1]$ com a medida de Lebesgue e defina $A_n=(0,1/n)$ e
$X_n=\mathds{1}_{A_n}$. Para cada $\omega\in(0,1]$, existe $n_0$ tal que
$1/n<\omega$ para todo $n\geqslant n_0$; para esses índices, $\omega\notin A_n$
e $X_n(\omega)=0$. Logo $X_n(\omega)\to0$ para todo $\omega$, e a convergência
é quase certa (na verdade, vale em todo ponto). Por outro lado, para
$0<\varepsilon\leqslant1$,
\[
\mathbb{P}\bigl(|X_n-0|\geqslant\varepsilon\bigr)=\mathbb{P}(A_n)=\frac{1}{n},
\]
de modo que $\sum_n\mathbb{P}(|X_n|\geqslant\varepsilon)=\infty$. A hipótese do
critério falha, mas a conclusão vale.
\end{example}

O que distingue o exemplo anterior de uma situação em que o critério é
necessário é a estrutura dos eventos. No exemplo, os eventos $A_n$ são
encaixados: $A_{n+1}\subseteq A_n$. Cada ponto sai de todos os $A_n$ a partir
de um certo índice, e a probabilidade de $A_n$ tende a zero sem que a soma
convirja. O primeiro lema de Borel-Cantelli não se aplica porque ele é um
critério baseado na soma, e a soma é grande justamente porque os eventos se
repetem nos mesmos pontos. Compare com o exemplo a seguir, no qual as
probabilidades são as mesmas, mas os eventos são independentes.

\begin{example}[Desvios independentes com série divergente]\label{ex:qc-independentes}
Seja $(A_n)$ uma sequência de eventos independentes com $\mathbb{P}(A_n)=1/n$ e
$X_n=\mathds{1}_{A_n}$. Pelo segundo lema de Borel-Cantelli,
$\mathbb{P}(\limsup_n A_n)=1$. Os eventos $A_n^c$ também são independentes e
$\sum_n\mathbb{P}(A_n^c)=\infty$, já que $\mathbb{P}(A_n^c)=1-1/n$ para
$n\geqslant2$; aplicando de novo o segundo lema,
$\mathbb{P}(\limsup_n A_n^c)=1$. Seja $\omega$ na interseção dos dois limites
superiores, que tem probabilidade $1$. Então $X_n(\omega)=1$ para infinitos
índices e $X_n(\omega)=0$ para infinitos índices, de modo que a sequência não
converge em $\omega$. Logo a probabilidade de convergência é zero.
\end{example}

Os dois exemplos têm a mesma distribuição marginal para cada $X_n$ e conduzem a
conclusões opostas. Isso mostra que a convergência quase certa não é uma
propriedade da coleção de distribuições individuais: ela depende da estrutura de
dependência entre as variáveis, isto é, da distribuição conjunta da sequência.
A convergência em distribuição, que estudaremos na \Cref{sec:distribuicao}, é
justamente a noção que ignora essa estrutura.

%----------------------------------------------------------------------
\section{Convergência em probabilidade}\label{sec:probabilidade}

\subsection{A definição}\label{subsec:prob-definicao}

A convergência em probabilidade é a noção que corresponde à pergunta
``qual é a chance de que $X_n$ se afaste de $X$ em uma etapa fixada?''. Ela não
tenta controlar o comportamento simultâneo de todos os índices, e é justamente
por isso que é mais fraca que a convergência quase certa.

\begin{definition}[Convergência em probabilidade]\label{def:prob}
Sejam $X,X_1,X_2,\ldots$ variáveis aleatórias em um mesmo espaço de
probabilidade. Dizemos que $(X_n)$ \emph{converge em probabilidade} para $X$ se,
para todo $\varepsilon>0$,
\[
\lim_{n\to\infty}\mathbb{P}\bigl(|X_n-X|\geqslant\varepsilon\bigr)=0.
\]
Nesse caso, escrevemos $X_n\xrightarrow{\mathbb{P}}X$.
\end{definition}

A definição é equivalente a dizer que, para todos $\varepsilon>0$ e
$\delta>0$, existe $n_0$ tal que
\[
\mathbb{P}\bigl(|X_n-X|<\varepsilon\bigr)\geqslant1-\delta
\qquad\text{para todo }n\geqslant n_0.
\]
Essa forma é útil quando queremos exibir $n_0$ explicitamente. A forma com o
limite é mais compacta e será a que usaremos nas demonstrações.

Duas observações sobre a definição. A primeira é que o evento cuja
probabilidade se controla depende de $\varepsilon$, mas não de $\omega$: a
exigência é que, em cada etapa, a probabilidade do desvio seja pequena, e não
que o desvio seja pequeno em cada ponto. A segunda é que a escolha de
$\geqslant$ ou $>$ no evento não altera a definição, porque
$\{|X_n-X|\geqslant\varepsilon\}\subseteq\{|X_n-X|>\varepsilon'\}$ para todo
$\varepsilon'<\varepsilon$. Usaremos $\geqslant$ para manter a coerência com as
desigualdades de Markov e Chebyshev, que são enunciadas nessa forma.

\subsection{Exemplos}\label{subsec:prob-exemplos}

\begin{example}[Convergência em probabilidade de uma translação]\label{ex:prob-translacao}
Seja $X$ uma variável aleatória e $X_n=X+1/n$. Para todo $\varepsilon>0$ e todo
$n>1/\varepsilon$,
\[
\mathbb{P}\bigl(|X_n-X|\geqslant\varepsilon\bigr)
=
\mathbb{P}\Bigl(\frac{1}{n}\geqslant\varepsilon\Bigr)=0,
\]
logo $X_n\to X$ em probabilidade. O exemplo é o mesmo da
\Cref{ex:qc-translacao}, e nele a convergência vale em todas as noções que
estudaremos.
\end{example}

\begin{example}[A frequência de um resultado no dado]\label{ex:prob-binomial}
Lancemos um dado honesto $n$ vezes e seja $N_n$ o número de vezes em que o
resultado é $5$ ou $6$. Então $N_n\sim\mathrm{Bin}(n,1/3)$, com
$\mathbb{E}[N_n]=n/3$ e $\operatorname{Var}(N_n)=2n/9$. Pela desigualdade de
Chebyshev aplicada a $N_n/n$,
\[
\mathbb{P}\Bigl(\Bigl|\frac{N_n}{n}-\frac{1}{3}\Bigr|\geqslant\varepsilon\Bigr)
\leqslant
\frac{1}{\varepsilon^2}\operatorname{Var}\Bigl(\frac{N_n}{n}\Bigr)
=
\frac{1}{\varepsilon^2}\cdot\frac{2n/9}{n^2}
=
\frac{2}{9n\varepsilon^2}\longrightarrow0.
\]
Logo $N_n/n\to1/3$ em probabilidade. O número $1/3$ é a probabilidade de o
resultado ser $5$ ou $6$ em um lançamento, e o exemplo é a versão mais simples
da lei dos grandes números.
\end{example}

\subsection{Convergência quase certa implica convergência em probabilidade}
\label{subsec:qc-implies-prob}

A convergência quase certa é mais forte que a convergência em probabilidade.
A prova é curta e usa a caracterização da \Cref{prop:qc-caracterizacao} junto
com a continuidade da medida.

\begin{proposition}\label{prop:qc-implies-prob}
Se $X_n\to X$ quase certamente, então $X_n\to X$ em probabilidade.
\end{proposition}

\begin{proof}
Fixe $\varepsilon>0$. Escolha $m$ com $1/m<\varepsilon$ e defina
$A_n=\{|X_n-X|\geqslant\varepsilon\}$ e $A_{n,m}=\{|X_n-X|\geqslant1/m\}$.
Como $A_n\subseteq A_{n,m}$, temos
$\limsup_n A_n\subseteq\limsup_n A_{n,m}$. Pela
\Cref{prop:qc-caracterizacao}, $\mathbb{P}(\limsup_n A_{n,m})=0$ e, portanto,
$\mathbb{P}(\limsup_n A_n)=0$. Considere a
sequência decrescente
\[
B_n=\bigcup_{k=n}^{\infty}A_k.
\]
Temos $A_n\subseteq B_n$ e $\bigcap_n B_n=\limsup_n A_n$. Pela continuidade da
medida, \eqref{eq:continuidade},
\[
\lim_{n\to\infty}\mathbb{P}(B_n)
=
\mathbb{P}\Bigl(\bigcap_{n=1}^{\infty}B_n\Bigr)
=
\mathbb{P}\Bigl(\limsup_{n\to\infty}A_n\Bigr)=0.
\]
Como $0\leqslant\mathbb{P}(A_n)\leqslant\mathbb{P}(B_n)$, segue que
$\mathbb{P}(A_n)\to0$, isto é, $X_n\to X$ em probabilidade.
\end{proof}

Vale sublinhar a estrutura da prova, porque ela se repete. Para passar de
``quase certamente'' para ``em probabilidade'', não comparamos $X_n$ com $X$
diretamente. Em vez disso, usamos que a não convergência quase certa é o evento
$\limsup_n A_n$, cuja probabilidade é zero, e que a probabilidade de $A_n$ é
majorada pela probabilidade da cauda $B_n$, que converge para a probabilidade
do limite superior. A continuidade da medida faz a ponte entre a informação
sobre o limite superior e a informação sobre cada termo.

\subsection{A recíproca é falsa: a sequência que varre o intervalo}
\label{subsec:typewriter}

A convergência em probabilidade não implica convergência quase certa. O
contraexemplo padrão é uma sequência de indicadoras de intervalos que percorrem
o intervalo $(0,1]$ em blocos de comprimento decrescente. Em cada etapa, a
probabilidade de o indicador valer $1$ é pequena, mas cada ponto é visitado
infinitas vezes.

\begin{example}[A sequência que varre o intervalo]\label{ex:typewriter}
Seja $\Omega=(0,1]$ com a medida de Lebesgue. Para cada $k\geqslant1$,
particione $(0,1]$ nos intervalos
\[
I_{k,j}=\Bigl(\frac{j-1}{k},\frac{j}{k}\Bigr],
\qquad j=1,\ldots,k,
\]
e enumere esses intervalos em blocos, pondo $I_{k,j}$ na posição
$n(k,j)=k(k-1)/2+j$. Defina
\[
X_n=\mathds{1}_{I_{k,j}}
\qquad\text{quando } n=n(k,j).
\]
Então $X_n\to0$ em probabilidade, mas $(X_n)$ não converge em nenhum ponto de
$\Omega$.
\end{example}

\begin{proof}[Verificação]
Para $n=n(k,j)$ temos $\mathbb{P}(X_n=1)=|I_{k,j}|=1/k$. Dado
$\varepsilon\in(0,1]$,
\[
\mathbb{P}\bigl(|X_n-0|\geqslant\varepsilon\bigr)
=
\mathbb{P}(X_n=1)=\frac{1}{k}\longrightarrow0
\qquad\text{quando } n\to\infty,
\]
porque $k\to\infty$ ao longo dos blocos. Logo $X_n\to0$ em probabilidade.

Fixe agora $\omega\in(0,1]$. Em cada bloco $k$, o ponto $\omega$ pertence a
exatamente um dos intervalos $I_{k,j}$, digamos $I_{k,j(k)}$, e não pertence aos
demais. No bloco $k$, portanto, $X_n(\omega)=1$ para um índice e
$X_n(\omega)=0$ para os outros $k-1$ índices. Como há infinitos blocos, a
sequência $X_n(\omega)$ vale $1$ em infinitos índices e $0$ em infinitos
índices. Uma sequência de números reais que assume os valores $0$ e $1$
infinitas vezes não converge. Logo $(X_n(\omega))$ não converge para nenhum
número real, e a convergência quase certa falha em todos os pontos.
\end{proof}

Vale entender por que o exemplo funciona, e não apenas verificar as contas. A
convergência em probabilidade olha para a etapa $n$ e pergunta qual é a massa
do conjunto onde $X_n$ é grande. No exemplo, essa massa é $1/k$, que tende a
zero. Mas os conjuntos onde $X_n$ é grande não param de se mover: cada ponto é
coberto uma vez em cada bloco. A convergência quase certa olha para cada ponto e
pergunta se o conjunto dos índices em que ele é coberto é finito. No exemplo,
esse conjunto é infinito para todos os pontos. As duas perguntas são diferentes,
e o exemplo as separa.

\begin{remark}[Subsequências]\label{rem:subsequencias}
Embora a convergência em probabilidade não implique convergência quase certa da
sequência inteira, ela implica a existência de uma subsequência que converge
quase certamente. Mais precisamente, $X_n\to X$ em probabilidade se, e somente
se, toda subsequência de $(X_n)$ possui uma subsubsequência que converge quase
certamente para $X$.
A implicação da esquerda para a direita se prova escolhendo índices $n_k$ com
$\mathbb{P}(|X_{n_k}-X|\geqslant1/k)\leqslant2^{-k}$ e aplicando a
\Cref{cor:criterio-bc-epsn}. A recíproca se prova por contraposição: se
$X_n\not\to X$ em probabilidade, existem $\varepsilon>0$, $\delta>0$ e uma
subsequência com $\mathbb{P}(|X_{n_k}-X|\geqslant\varepsilon)\geqslant\delta$
para todo $k$, e nenhuma subsubsequência dessa subsequência pode convergir quase
certamente, já que a convergência quase certa implica a convergência em
probabilidade pela \Cref{prop:qc-implies-prob}. O exemplo da
\Cref{ex:typewriter} mostra que a subsequência não pode ser substituída pela
sequência inteira.
\end{remark}

%----------------------------------------------------------------------
\section{Convergência em média e em média quadrática}\label{sec:medias}

\subsection{As definições}\label{subsec:medias-definicoes}

As noções desta seção substituem a pergunta ``o erro é pequeno?'' pela pergunta
``o erro médio é pequeno?''. Elas têm a vantagem de se reduzirem a cálculos de
esperança, que já sabemos fazer, e a desvantagem de serem sensíveis a poucos
pontos onde o erro é grande, como veremos nos exemplos.

\begin{definition}[Convergência em média]\label{def:media}
Sejam $X,X_1,X_2,\ldots$ variáveis aleatórias integráveis, isto é, com
$\mathbb{E}|X_n|<\infty$ para todo $n$ e $\mathbb{E}|X|<\infty$. Dizemos que
$(X_n)$ \emph{converge em média} para $X$ se
\[
\lim_{n\to\infty}\mathbb{E}\bigl[|X_n-X|\bigr]=0.
\]
Nesse caso, escrevemos $X_n\xrightarrow{L^1}X$.
\end{definition}

\begin{definition}[Convergência em média quadrática]\label{def:media-quadratica}
Sejam $X,X_1,X_2,\ldots$ variáveis aleatórias com $\mathbb{E}[X_n^2]<\infty$ e
$\mathbb{E}[X^2]<\infty$. Dizemos que $(X_n)$ \emph{converge em média
quadrática} para $X$ se
\[
\lim_{n\to\infty}\mathbb{E}\bigl[(X_n-X)^2\bigr]=0.
\]
Nesse caso, escrevemos $X_n\xrightarrow{L^2}X$.
\end{definition}

As definições pedem que os momentos envolvidos sejam finitos, pois do contrário
as esperanças não estariam definidas. No caso da média quadrática, a
desigualdade de Cauchy-Schwarz garante que $\mathbb{E}|X_n-X|$ também é finito,
de modo que a convergência em média quadrática só faz sentido quando a
convergência em média também faz.

\begin{example}[Convergência em média quadrática]\label{ex:media-quadratica}
Seja $Z_n$ uma variável aleatória com
\[
\mathbb{P}(Z_n=1)=\frac{1}{n},
\qquad
\mathbb{P}(Z_n=2)=1-\frac{1}{n}.
\]
Então $Z_n\to2$ em média quadrática, pois
\[
\mathbb{E}\bigl[(Z_n-2)^2\bigr]
=
(1-2)^2\cdot\frac{1}{n}+(2-2)^2\cdot\Bigl(1-\frac{1}{n}\Bigr)
=
\frac{1}{n}\longrightarrow0.
\]
O exemplo é o mesmo que aparece em Grimmett e Welsh e mostra que a convergência
em média quadrática pode ser verificada por um cálculo direto, sem invocar
nenhum teorema.
\end{example}

\begin{example}[A fronteira da média quadrática]\label{ex:media-quadratica-alfa}
Seja $\alpha>0$ e defina $Z_n$ por
\[
\mathbb{P}(Z_n=n^\alpha)=\frac{1}{n},
\qquad
\mathbb{P}(Z_n=0)=1-\frac{1}{n}.
\]
Então
\[
\mathbb{E}\bigl[Z_n^2\bigr]=n^{2\alpha}\cdot\frac{1}{n}=n^{2\alpha-1},
\]
de modo que $Z_n\to0$ em média quadrática se, e somente se, $2\alpha-1<0$, isto
é, $\alpha<1/2$. Por outro lado, para todo $\alpha>0$ e todo $\varepsilon>0$,
quando $n^\alpha\geqslant\varepsilon$ temos
\[
\mathbb{P}\bigl(|Z_n|\geqslant\varepsilon\bigr)
=
\mathbb{P}(Z_n=n^\alpha)
=
\frac{1}{n}\longrightarrow0,
\]
logo $Z_n\to0$ em probabilidade para todo $\alpha>0$. Assim, para
$1/2\leqslant\alpha<1$, a sequência converge em probabilidade e em média, mas
não em média quadrática. O exemplo será retomado na
\Cref{subsec:medias-contraexemplos}.
\end{example}

\subsection{As implicações}\label{subsec:medias-implicacoes}

A proposição seguinte reúne as comparações entre as convergências em média, em
média quadrática e em probabilidade. A parte relativa à distribuição será
completada na \Cref{sec:distribuicao}, depois de introduzirmos essa noção.

\begin{proposition}\label{prop:medias}
Sejam $X,X_1,X_2,\ldots$ variáveis aleatórias.
\begin{enumerate}
\item[(a)] Se $X_n\to X$ em média quadrática, então $X_n\to X$ em média.
\item[(b)] Se $X_n\to X$ em média quadrática, então
$X_n\to X$ em probabilidade.
\item[(c)] Se $X_n\to X$ em média, então $X_n\to X$ em probabilidade.
\end{enumerate}
\end{proposition}

\begin{proof}
Para (a), aplicamos a desigualdade de Cauchy-Schwarz à variável $Z=X_n-X$:
\[
\mathbb{E}|Z|
\leqslant
\bigl(\mathbb{E}[Z^2]\bigr)^{1/2}.
\]
Se $\mathbb{E}[(X_n-X)^2]\to0$, então $\mathbb{E}|X_n-X|\to0$.

Para (b), aplicamos a desigualdade de Chebyshev à mesma variável. Para
$\varepsilon>0$,
\[
\mathbb{P}\bigl(|X_n-X|\geqslant\varepsilon\bigr)
\leqslant
\frac{1}{\varepsilon^2}\mathbb{E}\bigl[(X_n-X)^2\bigr].
\]
Se o lado direito tende a zero, o lado esquerdo também tende a zero, e a
convergência em probabilidade segue.

Para (c), aplicamos a desigualdade de Markov a $|X_n-X|$:
\[
\mathbb{P}\bigl(|X_n-X|\geqslant\varepsilon\bigr)
\leqslant
\frac{1}{\varepsilon}\mathbb{E}\bigl[|X_n-X|\bigr],
\]
e o argumento é o mesmo de (b).
\end{proof}

A proposição mostra a cadeia
\[
L^2\ \Longrightarrow\ L^1\ \Longrightarrow\ \mathbb{P}.
\]
A primeira implicação é estrita, como mostra o
\Cref{ex:media-quadratica-alfa} com $1/2\leqslant\alpha<1$; a segunda também é
estrita, como mostra o exemplo a seguir.

\begin{remark}[Convergência em $L^p$]\label{rem:lp}
Poderíamos definir a convergência em média de ordem $p\geqslant1$ exigindo
$\mathbb{E}|X_n-X|^p\to0$. A desigualdade de Lyapunov, demonstrada na nota
\emph{Desigualdades de Markov e Chebyshev e o lema de Borel-Cantelli}, diz que
\[
\bigl(\mathbb{E}|Z|^q\bigr)^{1/q}
\leqslant
\bigl(\mathbb{E}|Z|^p\bigr)^{1/p}
\qquad\text{para }1\leqslant q\leqslant p,
\]
de modo que a convergência em $L^p$ implica a convergência em $L^q$ para todo
$q\leqslant p$. A implicação (a) da \Cref{prop:medias} é o caso $p=2$, $q=1$.
Nesta nota, tratamos apenas $p=1$ e $p=2$, que são os casos usados no curso; a
mesma proposição vale, com a mesma prova, para qualquer $p\geqslant1$.
\end{remark}

\begin{example}[Convergência em probabilidade sem convergência em média]
\label{ex:media-sem-quadratica}
Seja $Z_n$ com $\mathbb{P}(Z_n=n)=1/n$ e $\mathbb{P}(Z_n=0)=1-1/n$. Então
$Z_n\to0$ em probabilidade, pois $\mathbb{P}(|Z_n|\geqslant\varepsilon)=1/n\to0$
para todo $\varepsilon\in(0,1]$. Além disso,
\[
\mathbb{E}\bigl[|Z_n|\bigr]=\frac{1}{n}\cdot n=1\longrightarrow1,
\]
logo $Z_n$ não converge para $0$ em média. Tampouco converge em média quadrática,
pois $\mathbb{E}[Z_n^2]=n\to\infty$. O exemplo é o caso $\alpha=1$ do
\Cref{ex:media-quadratica-alfa}, e mostra que a massa $1/n$ sobre o valor $n$ é
pequena demais para afetar a probabilidade e grande demais para afetar a média.
\end{example}

\subsection{O que a média não vê}\label{subsec:medias-contraexemplos}

A convergência quase certa e a convergência em média não são comparáveis. Já
sabemos que a convergência em média implica convergência em probabilidade e,
portanto, nada impede que ela coexista com a convergência quase certa. Mas a
recíproca falha, e o exemplo a seguir é o mais simples possível.

\begin{example}[Convergência quase certa sem convergência em média]
\label{ex:qc-sem-media}
Seja $\Omega=(0,1]$ com a medida de Lebesgue e defina
\[
X_n=n\,\mathds{1}_{(0,1/n)}.
\]
Para cada $\omega\in(0,1]$, existe $n_0$ tal que $1/n<\omega$ para todo
$n\geqslant n_0$; para esses índices, $\omega\notin(0,1/n)$ e $X_n(\omega)=0$.
Logo $X_n(\omega)\to0$ para todo $\omega$, e a convergência é quase certa. Por
outro lado,
\[
\mathbb{E}|X_n|=n\cdot\frac{1}{n}=1
\qquad\text{para todo }n,
\]
de modo que $X_n$ não converge para $0$ em média. Também não converge em média
quadrática, pois $\mathbb{E}[X_n^2]=n\to\infty$. O exemplo mostra que a
convergência quase certa não controla a magnitude do erro nos pontos onde ele é
grande: em cada etapa, o erro vale $n$ em um conjunto de medida $1/n$.
\end{example}

O exemplo anterior é a razão pela qual o teorema da convergência dominada, que
estudaremos na \Cref{sec:quadro}, precisa de uma hipótese adicional. A
convergência quase certa garante que o erro tende a zero em cada ponto, mas não
garante que o erro médio tende a zero, porque a massa de pontos com erro grande
pode se concentrar em conjuntos cada vez menores sem que a integral diminua.

\begin{example}[Convergência em média sem convergência quase certa]
\label{ex:media-sem-qc}
A sequência da \Cref{ex:typewriter} converge em probabilidade e também em
média, pois
\[
\mathbb{E}|X_n|=\mathbb{P}(X_n=1)=\frac{1}{k}\longrightarrow0
\]
ao longo dos blocos. No entanto, ela não converge em nenhum ponto. Logo a
convergência em média não implica convergência quase certa.
\end{example}

Reunindo os exemplos, obtemos o quadro seguinte, que será completado na
\Cref{sec:distribuicao}.

\begin{center}
\begin{tabular}{@{}lll@{}}
\toprule
Sequência & Converge em & Não converge em \\
\midrule
$Z_n$ do \Cref{ex:media-quadratica-alfa}, $\alpha=0{,}7$ & $\mathbb{P}$, $L^1$ & $L^2$ \\
$n\mathds{1}_{(0,1/n)}$ & q.c., $\mathbb{P}$ & $L^1$, $L^2$ \\
$X_n$ do \Cref{ex:typewriter} & $\mathbb{P}$, $L^1$ & q.c. \\
\bottomrule
\end{tabular}
\end{center}

A tabela não esgota as possibilidades, mas cobre as separações que aparecem com
mais frequência. Note que nenhuma das linhas exibe convergência em média
quadrática sem convergência em média ou em probabilidade, o que é coerente com
a \Cref{prop:medias}.

\subsection{Um cálculo que será reutilizado}\label{subsec:medias-variancia}

Terminamos a seção com o cálculo que dá origem à lei fraca dos grandes números.
Sejam $X_1,X_2,\ldots$ variáveis aleatórias independentes com média $\mu$ e
variância $\sigma^2<\infty$, e seja $S_n=X_1+\cdots+X_n$. Como a variância de
uma soma de variáveis independentes é a soma das variâncias,
\[
\operatorname{Var}(S_n)=n\sigma^2.
\]
Logo
\[
\mathbb{E}\Bigl[\Bigl(\frac{S_n}{n}-\mu\Bigr)^2\Bigr]
=
\operatorname{Var}\Bigl(\frac{S_n}{n}\Bigr)
=
\frac{1}{n^2}\operatorname{Var}(S_n)
=
\frac{\sigma^2}{n}\longrightarrow0.
\]
Pela \Cref{prop:medias}, $S_n/n\to\mu$ em média quadrática, em média e em
probabilidade. Esse é o conteúdo da lei fraca dos grandes números na hipótese de
variância finita. A pergunta que a lei forte responde é se a convergência vale
também quase certamente, e a resposta exige um argumento diferente, porque a
cota $\sigma^2/n$ não é somável em $n$.

%----------------------------------------------------------------------
\section{Convergência em distribuição}\label{sec:distribuicao}

\subsection{A definição e os pontos de continuidade}\label{subsec:dist-definicao}

As noções anteriores comparam os valores de $X_n$ com os de $X$ e, por isso,
exigem que as variáveis vivam em um mesmo espaço de probabilidade. A
convergência em distribuição abandona essa comparação e olha apenas para as
funções de distribuição. Ela é a noção que aparece no teorema central do
limite, quando se diz que uma soma normalizada tem distribuição aproximadamente
normal sem que se especifique uma variável limite no mesmo espaço.

\begin{definition}[Convergência em distribuição]\label{def:dist}
Sejam $X,X_1,X_2,\ldots$ variáveis aleatórias com funções de distribuição
$F_n(x)=\mathbb{P}(X_n\leqslant x)$ e $F(x)=\mathbb{P}(X\leqslant x)$. Dizemos
que $(X_n)$ \emph{converge em distribuição} para $X$, e escrevemos
$X_n\Rightarrow X$, se
\[
\lim_{n\to\infty}F_n(x)=F(x)
\]
para todo $x\in\mathbb{R}$ em que $F$ é contínua.
\end{definition}

A restrição aos pontos de continuidade de $F$ não é um detalhe técnico
dispensável. Ela é o que torna a definição estável e o que permite enunciar os
teoremas de convergência sem hipóteses artificiais sobre as variáveis. O
exemplo seguinte mostra o que aconteceria se exigíssemos convergência em todos
os pontos.

\begin{example}[Por que os pontos de continuidade aparecem]\label{ex:dist-continuidade}
Seja $X\equiv0$ e $X_n=1/n$. A sequência converge para $X$ em todos os
sentidos estudados até aqui: quase certamente, em probabilidade, em média, em
média quadrática. As funções de distribuição são
$F(x)=\mathds{1}_{[0,\infty)}(x)$ e $F_n(x)=\mathds{1}_{[1/n,\infty)}(x)$.
Para $x<0$ e para $x>0$ com $n$ grande, $F_n(x)=F(x)$; mas em $x=0$ temos
$F_n(0)=0$ para todo $n$, enquanto $F(0)=1$. A convergência pontual das funções
de distribuição produz, no limite, a função $\mathds{1}_{(0,\infty)}$, que não é
contínua à direita em $0$ e não é a função de distribuição de $X$. A exigência
de convergência apenas nos pontos de continuidade de $F$ elimina o problema,
porque nesses pontos $F$ não tem salto e o valor pontual é determinado pelos
valores vizinhos. Note que $0$ é justamente o ponto de descontinuidade de $F$.
\end{example}

O exemplo anterior também explica por que a convergência em distribuição é uma
propriedade da \emph{distribuição} e não das variáveis. Se $X_n\Rightarrow X$ e
$Y_n$ tem a mesma distribuição que $X_n$ para todo $n$, então
$Y_n\Rightarrow X$. A relação entre $X_n$ e $X$ não desempenha nenhum papel,
apenas as funções de distribuição. Em particular, não é necessário que as
variáveis estejam definidas no mesmo espaço de probabilidade.

\subsection{Exemplos}\label{subsec:dist-exemplos}

\begin{example}[Discretização da uniforme]\label{ex:dist-discretizacao}
Para cada $n\geqslant1$, seja $X_n$ uniformemente distribuída no conjunto
$\{0,1/n,2/n,\ldots,(n-1)/n\}$ e seja $U$ uniforme em $(0,1)$. Então
$X_n\Rightarrow U$. De fato, para $x\in[0,1)$,
\[
F_n(x)=\mathbb{P}(X_n\leqslant x)
=
\frac{\#\{k\in\{0,\ldots,n-1\}:k/n\leqslant x\}}{n}
=
\frac{\lfloor nx\rfloor+1}{n}\longrightarrow x=F_U(x),
\]
e para $x<0$ ou $x\geqslant1$ as duas funções valem $0$ ou $1$ para todo $n$
grande. Como $F_U$ é contínua em todos os pontos, a convergência vale em todo
$x$. O exemplo é a versão determinística do fato de que uma amostra grande de
uma distribuição contínua a aproxima, e é o modelo discreto da integração de
Monte Carlo.
\end{example}

\begin{example}[Distribuição sem convergência em probabilidade]
\label{ex:dist-sem-prob}
Seja $U$ uma variável aleatória com $\mathbb{P}(U=-1)=\mathbb{P}(U=1)=1/2$ e
defina
\[
Z_n=
\begin{cases}
U, & n\text{ ímpar},\\
-U, & n\text{ par}.
\end{cases}
\]
Cada $Z_n$ tem a mesma distribuição de $U$, logo $Z_n\Rightarrow U$. Por outro
lado,
\[
Z_n-U=
\begin{cases}
0, & n\text{ ímpar},\\
-2U, & n\text{ par},
\end{cases}
\]
de modo que, para todo $m$,
\[
\mathbb{P}\bigl(|Z_{2m}-U|>1\bigr)=\mathbb{P}\Bigl(|U|>\frac{1}{2}\Bigr)=1.
\]
Assim $Z_n$ não converge para $U$ em probabilidade. O exemplo mostra que a
convergência em distribuição não controla o erro $|Z_n-U|$: as duas variáveis
têm sempre a mesma distribuição, mas em metade dos índices elas têm sinais
opostos.
\end{example}

\subsection{Convergência em probabilidade implica convergência em distribuição}
\label{subsec:prob-implies-dist}

O resultado seguinte é o análogo, para a convergência em distribuição, da
\Cref{prop:qc-implies-prob}. A prova usa um argumento de sanduíche que será
reutilizado em outras situações.

\begin{proposition}\label{prop:prob-implies-dist}
Se $X_n\to X$ em probabilidade, então $X_n\Rightarrow X$.
\end{proposition}

\begin{proof}
Seja $x$ um ponto de continuidade de $F$. Para $\varepsilon>0$, considere os
eventos
\[
\{X\leqslant x-\varepsilon\},
\qquad
\{X_n\leqslant x\},
\qquad
\{X\leqslant x+\varepsilon\}.
\]
Se $X\leqslant x-\varepsilon$ e $|X_n-X|<\varepsilon$, então $X_n\leqslant x$.
Logo
\[
\{X\leqslant x-\varepsilon\}\cap\{|X_n-X|<\varepsilon\}
\subseteq
\{X_n\leqslant x\}.
\]
Tomando complementares e usando a subaditividade,
\[
\mathbb{P}(X_n\leqslant x)
\geqslant
\mathbb{P}(X\leqslant x-\varepsilon)-\mathbb{P}(|X_n-X|\geqslant\varepsilon).
\]
Do mesmo modo, se $X_n\leqslant x$ e $|X_n-X|<\varepsilon$, então
$X<X_n+\varepsilon\leqslant x+\varepsilon$. Logo
\[
\{X_n\leqslant x\}\cap\{|X_n-X|<\varepsilon\}
\subseteq
\{X\leqslant x+\varepsilon\},
\]
e portanto
\[
\mathbb{P}(X_n\leqslant x)
\leqslant
\mathbb{P}(X\leqslant x+\varepsilon)+\mathbb{P}(|X_n-X|\geqslant\varepsilon).
\]
Reunindo as duas desigualdades,
\[
F(x-\varepsilon)-\mathbb{P}(|X_n-X|\geqslant\varepsilon)
\leqslant
F_n(x)
\leqslant
F(x+\varepsilon)+\mathbb{P}(|X_n-X|\geqslant\varepsilon).
\]
Fazendo $n\to\infty$ e usando a hipótese de convergência em probabilidade,
\[
F(x-\varepsilon)
\leqslant
\liminf_{n\to\infty}F_n(x)
\leqslant
\limsup_{n\to\infty}F_n(x)
\leqslant
F(x+\varepsilon).
\]
Como $x$ é ponto de continuidade de $F$, temos $F(x-\varepsilon)\to F(x)$
quando $\varepsilon\downarrow0$ e $F(x+\varepsilon)\to F(x)$ pela continuidade
à direita. Logo os limites inferior e superior coincidem com $F(x)$, isto é,
$F_n(x)\to F(x)$.
\end{proof}

\begin{corollary}[A convergência em média quadrática implica as demais]
\label{cor:l2-dist}
Se $X_n\to X$ em média quadrática, em média, em probabilidade ou quase
certamente, então $X_n\Rightarrow X$.
\end{corollary}

\begin{proof}
A convergência em média quadrática implica a convergência em média e a
convergência em probabilidade pela \Cref{prop:medias}; a convergência em média
implica a convergência em probabilidade pela mesma proposição; a convergência
quase certa implica a convergência em probabilidade pela
\Cref{prop:qc-implies-prob}. Em todos os casos, a
\Cref{prop:prob-implies-dist} fornece a convergência em distribuição.
\end{proof}

A proposição completa a cadeia de implicações
\[
L^2\ \Longrightarrow\ L^1\ \Longrightarrow\ \mathbb{P}
\ \Longrightarrow\ \text{distribuição},
\]
que liga as cinco noções estudadas. A
\Cref{prop:qc-implies-prob} acrescenta a convergência quase certa no início da
cadeia, mas apenas até a convergência em probabilidade:
\[
\text{q.c.}\ \Longrightarrow\ \mathbb{P}\ \Longrightarrow\ \text{distribuição}.
\]

\subsection{Quando a recíproca vale: o limite constante}\label{subsec:dist-constante}

A recíproca da \Cref{prop:prob-implies-dist} é falsa, como mostra o
\Cref{ex:dist-sem-prob}. Há, no entanto, um caso importante em que ela vale: se
o limite é uma constante, então convergência em distribuição e convergência em
probabilidade são equivalentes.

\begin{proposition}\label{prop:dist-constante}
Seja $c\in\mathbb{R}$. Então $X_n\Rightarrow c$ se, e somente se,
$X_n\to c$ em probabilidade.
\end{proposition}

\begin{proof}
A implicação de probabilidade para distribuição é a
\Cref{prop:prob-implies-dist}. Suponha agora que $X_n\Rightarrow c$ e seja
$\varepsilon>0$. A função de distribuição de $c$ é
$F(x)=\mathds{1}_{[c,\infty)}(x)$, que é contínua em todo $x\neq c$. Logo
$c-\varepsilon$ e $c+\varepsilon/2$ são pontos de continuidade de $F$ e
\[
\mathbb{P}\bigl(|X_n-c|\geqslant\varepsilon\bigr)
\leqslant
\mathbb{P}(X_n\leqslant c-\varepsilon)+\mathbb{P}\Bigl(X_n>c+\frac{\varepsilon}{2}\Bigr)
\leqslant
F_n(c-\varepsilon)+1-F_n\Bigl(c+\frac{\varepsilon}{2}\Bigr).
\]
Pela convergência em distribuição, o lado direito tende a
$F(c-\varepsilon)+1-F(c+\varepsilon/2)=0+1-1=0$. Isto é, $X_n\to c$ em
probabilidade.
\end{proof}

O resultado é útil porque, em muitas aplicações, o limite é uma constante
desconhecida, como a média de uma população ou uma probabilidade, e a conclusão
relevante é que o estimador se aproxima dela. A proposição diz que, nesse caso,
não há perda em trabalhar com convergência em distribuição.

\subsection{O que a convergência em distribuição não vê}\label{subsec:dist-oque-nao}

Vale resumir o que a convergência em distribuição controla e o que ela ignora.
Ela controla a massa que cada variável coloca em intervalos da forma
$(-\infty,x]$, ou seja, a forma das distribuições. Ela não controla a
dependência entre $X_n$ e $X$, como mostra o \Cref{ex:dist-sem-prob}; não
controla os momentos, no sentido de que $X_n\Rightarrow X$ não implica
$\mathbb{E}[X_n]\to\mathbb{E}[X]$; e não controla a probabilidade de eventos que
dependem conjuntamente de $X_n$ e $X$.

\begin{example}[Convergência em distribuição sem convergência de médias]
\label{ex:dist-momentos}
Seja $X\equiv0$ e $X_n=n\mathds{1}_{(0,1/n)}$, com $\Omega=(0,1]$ e medida de
Lebesgue. Para todo $x<0$, $F_n(x)=0=F(x)$. Para $x\geqslant0$ e $n>x$,
temos $F_n(x)=\mathbb{P}(X_n=0)=1-1/n$, que tende a $1=F(x)$; em $x=0$ o mesmo
cálculo se aplica. Logo $X_n\Rightarrow0$. No entanto,
$\mathbb{E}[X_n]=1$ para todo $n$, enquanto $\mathbb{E}[X]=0$. A convergência em
distribuição não vê a massa que escapa para valores grandes em conjuntos
pequenos.
\end{example}

O exemplo anterior é o mesmo da \Cref{ex:qc-sem-media}, lido sob outra ótica. A
sequência converge em todos os sentidos que não envolvem a magnitude do erro e
falha em todos os que a envolvem. A lição é que cada noção de convergência
responde a uma pergunta diferente, e a escolha da noção deve ser feita de acordo
com a pergunta.

%----------------------------------------------------------------------
\section{O quadro geral das implicações}\label{sec:quadro}

\subsection{O diagrama}\label{subsec:diagrama}

Reunimos na \Cref{fig:implicacoes} as implicações demonstradas até aqui. As
setas contínuas valem sem hipóteses adicionais; a seta tracejada representa o
teorema da convergência dominada, que será enunciado na
\Cref{subsec:dominada} e que exige uma cota integrável para o módulo das
variáveis.

\begin{figure}[htbp]
\centering
\begin{tikzpicture}[
  node distance=1.6cm,
  caixa/.style={draw, rounded corners=3pt, fill=white, inner sep=6pt,
    minimum width=1.5cm, font=\small},
  seta/.style={-{Latex[length=2.2mm]}, semithick},
  rotulo/.style={font=\scriptsize, fill=white, inner sep=2pt}
]
\node[caixa] (L2) at (0,2.8) {$L^2$};
\node[caixa] (L1) at (3.6,2.8) {$L^1$};
\node[caixa] (P)  at (2.4,0.2) {$\mathbb{P}$};
\node[caixa] (D)  at (6.0,0.2) {distribuição};
\node[caixa] (QC) at (-1.9,1.4) {q.c.};

\draw[seta] (L2) -- node[rotulo, above] {Cauchy--Schwarz} (L1);
\draw[seta] (L1) -- node[rotulo, right, pos=0.5] {Markov} (P);
\draw[seta] (QC) -- node[rotulo, below, pos=0.55] {\Cref{prop:qc-implies-prob}} (P);
\draw[seta] (P) -- node[rotulo, above] {\Cref{prop:prob-implies-dist}} (D);
\draw[seta, dashed] (QC) --
  node[rotulo, below, pos=0.3] {com dominação} (L1.south west);
\end{tikzpicture}
\caption{Implicações entre os modos de convergência. As setas contínuas valem
em geral. A seta tracejada representa o teorema da convergência dominada: a
convergência quase certa, acompanhada de uma dominância integrável, implica a
convergência em média e, por intermédio dela, todas as demais.}
\label{fig:implicacoes}
\end{figure}

A leitura do diagrama deve ser feita com cuidado. O fato de existir uma seta de
$L^2$ para $L^1$ não significa que toda sequência que converge em média
quadrática converge em média por um argumento diferente; significa que a
implicação vale e que a prova passa por Cauchy-Schwarz. Do mesmo modo, a seta
de $\mathbb{P}$ para distribuição tem uma prova própria, que não passa por
média alguma. O diagrama é um resumo das relações lógicas, não um roteiro de
demonstração.

\subsection{A tabela}\label{subsec:tabela}

A \Cref{tab:modos} reúne as definições, as implicações e, em cada linha, um
exemplo que refuta uma das recíprocas.

\begin{table}[htbp]
\centering
\footnotesize
\setlength{\tabcolsep}{4pt}
\begin{tabular}{@{}p{2.4cm}p{4.1cm}p{2.2cm}p{5.1cm}@{}}
\toprule
Modo & Definição & Implica & Exemplo que separa \\
\midrule
q.c. &
$\mathbb{P}(X_n\to X)=1$ &
$\mathbb{P}$, distribuição &
\Cref{ex:typewriter}: converge em $\mathbb{P}$, não em q.c. \\
em probabilidade &
$\mathbb{P}(|X_n-X|\geqslant\varepsilon)\to0$ &
distribuição &
\Cref{ex:dist-sem-prob}: converge em distribuição, não em
$\mathbb{P}$ \\
em média &
$\mathbb{E}|X_n-X|\to0$ &
$\mathbb{P}$, distribuição &
\Cref{ex:qc-sem-media}: converge q.c., não em média \\
em média quadrática &
$\mathbb{E}[(X_n-X)^2]\to0$ &
média, $\mathbb{P}$, distribuição &
\Cref{ex:media-quadratica-alfa} com $\alpha=0{,}7$ \\
em distribuição &
$F_n(x)\to F(x)$ nos pontos de continuidade &
nenhuma das anteriores &
\Cref{ex:dist-momentos}: as médias não convergem \\
\bottomrule
\end{tabular}
\caption{Resumo dos cinco modos de convergência. A coluna ``implica'' lista
apenas as implicações que valem sem hipóteses adicionais. A última coluna
indica um exemplo, neste texto, que ilustra a separação entre a noção da linha
e as demais: ele mostra uma implicação recíproca que falha.}
\label{tab:modos}
\end{table}

Duas observações sobre a tabela. A primeira é que a coluna ``implica'' não é
transitiva por linhas: a convergência em média quadrática implica a convergência
em distribuição, mas a tabela registra apenas a implicação para a média e para
a probabilidade, porque a implicação para a distribuição se obtém compondo com
a \Cref{prop:prob-implies-dist}. A segunda é que cada exemplo da última coluna
refuta a recíproca de alguma implicação que envolve a noção da linha, e não
pretende esgotar as separações possíveis.

\subsection{Convergência dominada: quando a convergência quase certa basta}
\label{subsec:dominada}

Falta tratar de uma situação que aparece com frequência nas aplicações. Muitas
vezes sabemos que $X_n\to X$ quase certamente e queremos concluir que as médias
também convergem. O \Cref{ex:qc-sem-media} mostra que isso não vale em geral: a
convergência quase certa não controla a magnitude do erro. O teorema da
convergência dominada dá a hipótese adicional que corrige o problema.

\begin{theorem}[Convergência dominada]\label{thm:dominada}
Sejam $X,X_1,X_2,\ldots$ variáveis aleatórias tais que $X_n\to X$ quase
certamente. Suponha que exista uma variável aleatória $Y$ com
$\mathbb{E}[Y]<\infty$ e
\[
|X_n|\leqslant Y
\qquad\text{quase certamente, para todo }n.
\]
Então $X$ é integrável, $\mathbb{E}|X_n-X|\to0$ e, em particular,
$\mathbb{E}[X_n]\to\mathbb{E}[X]$.
\end{theorem}

\begin{proof}
Como $X_n\to X$ quase certamente, $|X_n-X|\to0$ quase certamente. Além disso,
\[
|X_n-X|\leqslant|X_n|+|X|\leqslant2Y
\qquad\text{quase certamente,}
\]
sendo a última desigualdade consequência de $|X|\leqslant Y$ quase certamente
(que se obtém passando ao limite em $|X_n|\leqslant Y$). A função $2Y$ é
integrável e domina a sequência $|X_n-X|$. Pelo teorema da convergência
dominada de Lebesgue (na forma para variáveis aleatórias, idêntica à enunciada
na nota de funções de distribuição),
\[
\lim_{n\to\infty}\mathbb{E}|X_n-X|=\mathbb{E}\bigl[\lim_{n\to\infty}|X_n-X|\bigr]=0.
\]
A conclusão sobre as esperanças segue de
$|\mathbb{E}[X_n]-\mathbb{E}[X]|\leqslant\mathbb{E}|X_n-X|$.
\end{proof}

\begin{corollary}\label{cor:dominada}
Nas hipóteses da \Cref{thm:dominada}, a sequência converge em média, em
probabilidade e em distribuição. Se, além disso, $\mathbb{E}[Y^2]<\infty$,
então a convergência também vale em média quadrática.
\end{corollary}

\begin{proof}
A convergência em média é a \Cref{thm:dominada}; a convergência em
probabilidade segue da \Cref{prop:medias}(c) e a convergência em distribuição
segue da \Cref{prop:prob-implies-dist}. Para a convergência em média quadrática,
observe que $(X_n-X)^2\leqslant4Y^2$, que é integrável, e aplique novamente o
teorema da convergência dominada.
\end{proof}

O corolário reúne as consequências da convergência dominada: quando a sequência
é dominada por uma variável integrável, a convergência quase certa implica
todas as demais noções. A hipótese de dominância é
essencial: sem ela, o \Cref{ex:qc-sem-media} mostra que a convergência quase
certa pode coexistir com a divergência das médias.

\begin{example}[O máximo de uniformes converge em todos os modos]
\label{ex:maximo-todos}
Sejam $U_1,U_2,\ldots$ independentes e uniformes em $(0,1)$ e
$M_n=\max\{U_1,\ldots,U_n\}$. Na \Cref{ex:qc-maximo}, mostramos que
$M_n\to1$ quase certamente. Como $0\leqslant M_n\leqslant1$ e $1$ é
integrável, a \Cref{thm:dominada} se aplica com $Y\equiv1$. Logo
$M_n\to1$ em média, em probabilidade e em distribuição. Como $Y^2=1$ também é
integrável, o \Cref{cor:dominada} garante a convergência em média quadrática.

Podemos verificar o resultado por cálculo direto. Para $t\in[0,1]$,
\[
\mathbb{P}(M_n\leqslant t)=t^n,
\]
logo a densidade de $M_n$ é $nt^{n-1}$ em $(0,1)$ e
\[
\mathbb{E}[M_n]=\int_0^1t\cdot nt^{n-1}\,dt=\frac{n}{n+1}\longrightarrow1.
\]
Do mesmo modo,
\[
\mathbb{E}[M_n^2]=\int_0^1t^2\cdot nt^{n-1}\,dt=\frac{n}{n+2}\longrightarrow1,
\]
e portanto
\[
\mathbb{E}\bigl[(1-M_n)^2\bigr]
=
1-2\,\mathbb{E}[M_n]+\mathbb{E}[M_n^2]
=
\frac{2}{(n+1)(n+2)}\longrightarrow0.
\]
O exemplo é instrutivo porque a convergência em todos os modos pode ser
verificada de duas maneiras independentes: pelo teorema da convergência
dominada, que não usa a forma da distribuição, e pelo cálculo explícito dos
momentos, que não usa o teorema.
\end{example}

\subsection{Leitura final}\label{subsec:leitura-final}

As cinco noções não formam uma hierarquia linear, mas formam uma rede com duas
entradas fortes, a convergência quase certa e a convergência em média
quadrática, e uma saída fraca, a convergência em distribuição. A convergência em
probabilidade é o ponto de encontro: ela é consequência de ambas as entradas e
é a hipótese que basta para a convergência em distribuição. A convergência em
distribuição é a única que não exige um espaço comum e a única que não controla
o erro.

Na prática, a escolha da noção é ditada pelo que se pode estimar. Se o problema
fornece cotas para $\mathbb{P}(|X_n-X|\geqslant\varepsilon)$, a conclusão
natural é a convergência em probabilidade. Se fornece cotas para
$\mathbb{E}[(X_n-X)^2]$, a conclusão é a convergência em média quadrática, que é
mais forte. Se fornece cotas somáveis em $n$, a conclusão é a convergência quase
certa, pelo \Cref{prop:criterio-bc}. E se o interesse está apenas na forma
assintótica da distribuição, como no teorema central do limite, a convergência
em distribuição é a noção apropriada.

A nota sobre a lei forte dos grandes números explora exatamente a passagem de
cotas não somáveis para cotas somáveis. Lá, a cota de Chebyshev para a média
empírica é da ordem de $1/n$, que tende a zero mas não é somável; elevando ao
quadrado o erro, ou truncando as variáveis, obtém-se uma cota de ordem $1/n^2$,
somável, e o \Cref{prop:criterio-bc} produz a convergência quase certa.

%----------------------------------------------------------------------
\section{Exercícios propostos}\label{sec:exercicios}

\begin{exercise}\label{exercise:qc-propriedades}
Sejam $X,X_1,X_2,\ldots$ variáveis aleatórias.
\begin{enumerate}
\item[(a)] Mostre que $X_n\to X$ quase certamente se, e somente se,
$X_n-X\to0$ quase certamente.
\item[(b)] Mostre que, se $X_n\to X$ quase certamente, então
$|X_n|\to|X|$ quase certamente.
\item[(c)] Mostre que, se $X_n\to X$ em probabilidade, então
$|X_n|\to|X|$ em probabilidade.
\end{enumerate}
\end{exercise}

\begin{exercise}\label{exercise:unicidade}
Suponha que $X_n\to X$ quase certamente e $X_n\to Y$ quase certamente. Mostre
que $X=Y$ quase certamente. A conclusão continua válida se a convergência for
apenas em probabilidade? E se for apenas em distribuição?
\end{exercise}

\begin{exercise}\label{exercise:soma-prob}
Sejam $X_n\to X$ e $Y_n\to Y$ em probabilidade. Mostre que
$X_n+Y_n\to X+Y$ em probabilidade. \emph{Sugestão:} use a desigualdade
$|X_n+Y_n-X-Y|\leqslant|X_n-X|+|Y_n-Y|$ e a subaditividade da probabilidade.
\end{exercise}

\begin{exercise}\label{exercise:dist-ncg}
Seja $X$ com distribuição normal padrão e defina $X_n=(-1)^nX$. Mostre que
$X_n\Rightarrow X$, mas $(X_n)$ não converge em probabilidade para $X$.
Compare com o \Cref{ex:dist-sem-prob}.
\end{exercise}

\begin{exercise}\label{exercise:deslocamento-dist}
Sejam $X,X_1,X_2,\ldots$ variáveis aleatórias com $X_n\Rightarrow X$ e seja
$(c_n)$ uma sequência de números reais com $c_n\to0$. Mostre que
$X_n+c_n\Rightarrow X$.
\end{exercise}

\begin{exercise}\label{exercise:media-variancia}
Suponha que $X_n\to X$ em média quadrática. Mostre que
$\mathbb{E}[X_n]\to\mathbb{E}[X]$ e
$\operatorname{Var}(X_n)\to\operatorname{Var}(X)$. Dê um exemplo em que
$\mathbb{E}[X_n]\to\mathbb{E}[X]$ e, no entanto, $X_n$ não converge para $X$ em
média quadrática.
\end{exercise}

\begin{exercise}\label{exercise:bc-somavel}
Seja $(A_n)$ uma sequência de eventos com $\mathbb{P}(A_n)=2^{-n}$. Mostre que
\[
\sum_{n=1}^{\infty}\mathds{1}_{A_n}(\omega)<\infty
\qquad\text{para quase todo }\omega.
\]
Interprete o resultado em termos do número de eventos $A_n$ que ocorrem.
\end{exercise}

\begin{exercise}\label{exercise:frequencia-digitos}
Seja $\Omega=\{0,1,\ldots,9\}^{\mathbb{N}}$, com a medida produto determinada
por $\mathbb{P}(X_n=k)=1/10$, onde $X_n(\omega)=\omega_n$. Para $k$ fixo,
defina
\[
F_{n,k}=\frac{1}{n}\sum_{j=1}^{n}\mathds{1}_{\{X_j=k\}}.
\]
\begin{enumerate}
\item[(a)] Calcule $\mathbb{E}[F_{n,k}]$ e $\operatorname{Var}(F_{n,k})$.
\item[(b)] Conclua que $F_{n,k}\to1/10$ em probabilidade.
\item[(c)] Mostre que a cota obtida em (b) não é somável em $n$ e explique por
que o critério de Borel-Cantelli não se aplica diretamente. A convergência
quase certa será demonstrada na nota sobre a lei forte dos grandes números.
\end{enumerate}
\end{exercise}

\begin{starredexercise}\label{exercise:metrica-prob}
No conjunto das classes de equivalência de variáveis aleatórias, identificando
$X$ e $Y$ quando $X=Y$ quase certamente, defina
\[
d(X,Y)=\mathbb{E}\bigl[\min\{1,|X-Y|\}\bigr].
\]
\begin{enumerate}
\item[(a)] Mostre que $d$ está bem definida, que $d(X,Y)=d(Y,X)$ e que
$d(X,Y)=0$ se, e somente se, $X=Y$ quase certamente.
\item[(b)] Mostre a desigualdade triangular.
\item[(c)] Mostre que $d(X_n,X)\to0$ se, e somente se, $X_n\to X$ em
probabilidade.
\end{enumerate}
\end{starredexercise}

\begin{starredexercise}\label{exercise:subsequencia-qc}
Suponha que $X_n\to X$ em probabilidade. Construa, indutivamente, uma
subsequência $(X_{n_k})$ tal que
\[
\mathbb{P}\bigl(|X_{n_k}-X|\geqslant\tfrac{1}{k}\bigr)\leqslant2^{-k}
\qquad\text{para todo }k\geqslant1,
\]
e conclua que $X_{n_k}\to X$ quase certamente. Use o resultado para provar a
implicação da esquerda para a direita no \Cref{rem:subsequencias}.
\end{starredexercise}

%----------------------------------------------------------------------
% BIBLIOGRAFIA
%----------------------------------------------------------------------
\bibliographystyle{plain}
\bibliography{../arquivos-aux/referencias}

\end{document}
